\documentclass[12pt,a4paper]{article}
\usepackage[utf8]{inputenc}
\usepackage{amsmath,amssymb,amsthm}
\usepackage{bm}
\usepackage{mathtools}
\usepackage{booktabs}
\usepackage{array}
\usepackage{multirow}
\usepackage{geometry}
\usepackage{hyperref}
\usepackage{algorithm2e}

\geometry{margin=2.5cm}

\DeclareMathOperator{\tr}{tr}
\DeclareMathOperator{\rank}{rank}
\DeclareMathOperator{\Real}{Re}
\DeclareMathOperator{\vect}{vec}
\newcommand{\mat}[1]{\mathbf{#1}}
\newcommand{\norm}[1]{\left\|#1\right\|}

\newtheorem{theorem}{Theorem}
\newtheorem{lemma}{Lemma}
\newtheorem{proposition}{Proposition}
\newtheorem{remark}{Remark}
\newtheorem{assumption}{Assumption}

\title{\textbf{Mathematical Foundation for Cell-Free ISAC}\\\large Physical Modeling and Convex Reformulation (Plates I \& II)}
\author{ISAC Research Team}
\date{\today}

\begin{document}
\maketitle
\tableofcontents
\newpage

\section{System Modeling and Problem Formulation (Plate I: Physical Modeling)}
%==============================================================================

\subsection{Optimization Problem (Standard Form)}

Consider a Cell-Free ISAC system with $M$ distributed APs, $K$ communication users, and $P$ sensing targets. Each AP is equipped with $N_t$ transmit antennas. The optimization objective is to minimize the total transmit power while satisfying communication QoS, sensing detection, tracking accuracy, and AP selection constraints.

\begin{align}
\min_{\substack{\{\vect{w}_{m,k}\}, \{\mat{Z}_m\}, \\ \{b_{mp}\}}} \quad & \sum_{m=1}^{M} \left( \sum_{k=1}^{K} \norm{\vect{w}_{m,k}}^2 + \tr(\mat{Z}_m) \right) \tag{P1} \\
\text{s.t.} \quad & \text{SINR}_k^{\text{wc}} \geq \gamma_k, \quad \forall k \in \mathcal{K} \label{eq:sinr_comm} \\
& \text{SINR}_{S,p} \geq \gamma_S^{\text{PoD}}, \quad \forall p \in \mathcal{P} \label{eq:sinr_sens} \\
& \tr(\mat{J}_p^{\text{data}}) \geq \Gamma_{\text{Track}, p}, \quad \forall p \in \mathcal{P} \label{eq:pcrb} \\
& \sum_{m=1}^{M} b_{mp} = N_{\text{req}}, \quad \forall p \in \mathcal{P}^{\text{active}} \label{eq:ap_select} \\
& \sum_{k=1}^{K} \norm{\vect{w}_{m,k}}^2 + \tr(\mat{Z}_m) \leq P_{\max}, \quad \forall m \in \mathcal{M} \label{eq:power} \\
& \mat{Z}_m \succeq \mat{0}, \quad \forall m \in \mathcal{M} \label{eq:psd} \\
& b_{mp} \in \{0, 1\}, \quad \forall m \in \mathcal{M}, p \in \mathcal{P} \label{eq:binary}
\end{align}

where $\vect{w}_{m,k} \in \mathbb{C}^{N_t}$ denotes the communication beamforming vector from AP $m$ to user $k$, $\mat{Z}_m \in \mathbb{C}^{N_t \times N_t}$ represents the sensing covariance matrix at AP $m$, and $b_{mp} \in \{0,1\}$ is the binary indicator variable for whether AP $m$ serves target $p$. The parameters $P_{\max}$, $\gamma_k$, $\gamma_S^{\text{PoD}}$, $\Gamma_{\text{Track},p}$, and $N_{\text{req}}$ denote the per-AP maximum transmit power, the SINR threshold for user $k$, the sensing detection probability threshold, the tracking accuracy threshold (FIM trace constraint) for target $p$, and the required number of serving APs per target, respectively.

\subsection{Notation Table}

Table~\ref{tab:symbols} summarizes the main notation used throughout this paper.

\begin{table}[htbp]
\centering
\caption{Main Notation Table}\label{tab:symbols}
\begin{tabular}{cll}
\toprule
Symbol & Dimension & Physical Meaning \\
\midrule
$M$ & Scalar & Total number of APs \\
$N_t$ & Scalar & Number of transmit antennas per AP \\
$K$ & Scalar & Total number of communication users \\
$P$ & Scalar & Total number of sensing targets \\
$P_{\max}$ & Scalar & Maximum transmit power per AP \\
$\gamma_k$ & Scalar & SINR threshold for user $k$ \\
$\gamma_S^{\text{PoD}}$ & Scalar & Sensing detection probability threshold \\
$\Gamma_{\text{Track},p}$ & Scalar & Tracking accuracy threshold for target $p$ \\
$\sigma_c^2$ & Scalar & Communication receiver noise variance \\
$\sigma_s^2$ & Scalar & Sensing receiver noise variance \\
$\epsilon_h$ & Scalar & Relative error bound for communication channel estimation \\
$\epsilon_g$ & Scalar & Relative error bound for sensing channel estimation \\
$N_{\text{req}}$ & Scalar & Required number of serving APs per target \\
$\vect{h}_{m,k}$ & $\mathbb{C}^{N_t}$ & Communication channel from AP $m$ to user $k$ \\
$\vect{g}_{m,p}$ & $\mathbb{C}^{N_t}$ & Sensing channel from AP $m$ to target $p$ \\
$\vect{w}_{m,k}$ & $\mathbb{C}^{N_t}$ & Communication beam from AP $m$ to user $k$ \\
$\mat{Z}_m$ & $\mathbb{C}^{N_t \times N_t}$ & Sensing covariance matrix at AP $m$ \\
$b_{mp}$ & $\{0,1\}$ & Whether AP $m$ serves target $p$ \\
\bottomrule
\end{tabular}
\end{table}

\begin{remark}
Specific numerical parameters (e.g., $M=16$, $N_t=4$, $P_{\max}=30$ W) will be provided in the Simulation Setup section. The theoretical derivations herein maintain general symbolic notation.
\end{remark}

%==============================================================================
\section{System Model and Imperfect CSI}
%==============================================================================

The second step of physical modeling: characterizing the transmission of communication and sensing signals under imperfect channel state information.

\subsection{Communication Signal Model}

The received signal at user $k$ can be expressed as:
\begin{equation}
y_k = \underbrace{\sum_{m=1}^M \vect{h}_{m,k}^H \vect{w}_{m,k} s_k}_{\text{desired signal}} + \underbrace{\sum_{j \neq k} \sum_{m=1}^M \vect{h}_{m,k}^H \vect{w}_{m,j} s_j}_{\text{multi-user interference}} + n_k
\end{equation}
where $n_k \sim \mathcal{CN}(0, \sigma_c^2)$ denotes the receiver noise.

To facilitate analysis, we introduce stacked vector forms. Define the global communication channel $\vect{h}_k \in \mathbb{C}^{MN_t}$ and global communication beam $\vect{w}_k \in \mathbb{C}^{MN_t}$ as:
\begin{equation}
\vect{h}_k = [\vect{h}_{1,k}^T, \vect{h}_{2,k}^T, \ldots, \vect{h}_{M,k}^T]^T, \quad
\vect{w}_k = [\vect{w}_{1,k}^T, \vect{w}_{2,k}^T, \ldots, \vect{w}_{M,k}^T]^T
\end{equation}

The received signal can then be compactly written as:
\begin{equation}
y_k = \vect{h}_k^H \vect{w}_k s_k + \sum_{j \neq k} \vect{h}_k^H \vect{w}_j s_j + n_k
\end{equation}

\subsection{Sensing Signal Model}

For monostatic radar sensing, the reflected signal from target $p$ is:
\begin{equation}
y_{S,p} = \sum_{m=1}^M \vect{g}_{m,p}^H \mat{Z}_m \vect{g}_{m,p} + n_{S,p}
\end{equation}
where $n_{S,p} \sim \mathcal{CN}(0, \sigma_s^2)$ denotes the sensing receiver noise.

\subsection{Imperfect CSI Model}

In practical systems, channel estimation suffers from errors. This paper adopts the bounded error model:

\textbf{Communication channel}:
\begin{equation}
\vect{h}_k = \hat{\vect{h}}_k + \Delta\vect{h}_k, \quad \norm{\Delta\vect{h}_k} \leq \epsilon_h \norm{\hat{\vect{h}}_k}
\end{equation}

\textbf{Sensing channel}:
\begin{equation}
\vect{g}_p = \hat{\vect{g}}_p + \Delta\vect{g}_p, \quad \norm{\Delta\vect{g}_p} \leq \epsilon_g \norm{\hat{\vect{g}}_p}
\end{equation}

where $\hat{\vect{h}}_k$ and $\hat{\vect{g}}_p$ are the estimated channels, $\Delta\vect{h}_k$ and $\Delta\vect{g}_p$ are the estimation errors, and $\epsilon_h$ and $\epsilon_g$ are the relative error bounds.

%==============================================================================
\section{Communication Constraint Derivation}
%==============================================================================

The third step of physical modeling: starting from the signal model, we derive the original non-convex expressions of worst-case SINR, sensing SINR, and PCRB.

\subsection{SINR Definition}

The signal-to-interference-plus-noise ratio (SINR) for communication user $k$ is defined as:
\begin{equation}
\text{SINR}_k = \frac{|\vect{h}_k^H \vect{w}_k|^2}{\sum_{j \neq k} |\vect{h}_k^H \vect{w}_j|^2 + \sigma_c^2} \tag{1}
\end{equation}

\subsection{Worst-Case SINR Analysis}

Under imperfect CSI, we need to guarantee that the worst-case SINR meets the threshold requirement. Considering the worst cases for both the numerator and interference:

\textbf{Worst-case numerator}: When the error direction opposes the beam, the effective channel gain is minimized:
\begin{equation}
|\vect{h}_k^H \vect{w}_k|^2 \approx |\hat{\vect{h}}_k^H \vect{w}_k|^2 (1 - \epsilon_h)^2
\end{equation}

\textbf{Worst-case interference}: When the error direction enhances interference:
\begin{equation}
|\vect{h}_k^H \vect{w}_j|^2 \approx |\hat{\vect{h}}_k^H \vect{w}_j|^2 (1 + \epsilon_h)^2
\end{equation}

Therefore, the worst-case SINR can be approximated as:
\begin{equation}
\text{SINR}_k^{\text{wc}} \approx \text{SINR}_k^{\text{nom}} \cdot \frac{(1-\epsilon_h)^2}{(1+\epsilon_h)^2} \tag{2}
\end{equation}

\subsection{Robustness Factor and Threshold}

Define the communication robustness factor:
\begin{equation}
\eta_h = \left(\frac{1-\epsilon_h}{1+\epsilon_h}\right)^2 \tag{3}
\end{equation}

For $\epsilon_h = 0.10$ (10\% error), $\eta_h \approx 0.669$ ($-1.75$ dB). This implies that CSI errors cause an effective SINR degradation of approximately 1.75 dB.

To ensure the worst-case SINR meets the threshold, the nominal SINR requirement must be raised:
\begin{equation}
\gamma_k^{\text{robust}} = \gamma_k \cdot \frac{(1+\epsilon_h)^2}{(1-\epsilon_h)^2} = \frac{\gamma_k}{\eta_h} \tag{4}
\end{equation}

For $\gamma_k = 1$ (0 dB) and $\epsilon_h = 0.10$: $\gamma_k^{\text{robust}} \approx 1.493$ ($1.74$ dB).

%==============================================================================
\section{Sensing Constraint Derivation}
%==============================================================================

\subsection{Sensing SINR}

The sensing SINR for target $p$ is:
\begin{equation}
\text{SINR}_{S,p} = \frac{|\vect{g}_p^H \vect{z}_p|^2}{\sigma_s^2 \norm{\vect{z}_p}^2} \tag{5}
\end{equation}

By the Cauchy-Schwarz inequality, the optimal sensing beam is the matched filter:
\begin{equation}
\vect{z}_p^* = \sqrt{P_{S,p}} \frac{\vect{g}_p}{\norm{\vect{g}_p}} \tag{6}
\end{equation}

yielding the optimal SINR:
\begin{equation}
\text{SINR}_{S,p}^* = \frac{P_{S,p} \norm{\vect{g}_p}^2}{\sigma_s^2} \tag{7}
\end{equation}

To meet the detection threshold $\gamma_S^{\text{PoD}}$, the minimum sensing power is:
\begin{equation}
P_{S,p}^{\min} = \frac{\gamma_S^{\text{PoD}} \sigma_s^2}{\norm{\vect{g}_p}^2} \tag{8}
\end{equation}

\subsection{Robust Sensing Constraint}

Similarly, define the sensing robustness factor:
\begin{equation}
\eta_g = \left(\frac{1-\epsilon_g}{1+\epsilon_g}\right)^2 \tag{9}
\end{equation}

For $\epsilon_g = 0.15$ (15\% error), $\eta_g \approx 0.546$ ($-2.63$ dB).

The robust sensing threshold:
\begin{equation}
(\gamma_S^{\text{PoD}})^{\text{robust}} = \gamma_S^{\text{PoD}} \cdot \frac{(1+\epsilon_g)^2}{(1-\epsilon_g)^2} = \frac{\gamma_S^{\text{PoD}}}{\eta_g} \tag{10}
\end{equation}

\subsection{PCRB Constraint and Fisher Information Matrix}

The tracking accuracy constraint is characterized through the trace of the Fisher Information Matrix (FIM). The data-dependent FIM is:
\begin{equation}
\mat{J}_p^{\text{data}} = \frac{2}{\sigma_s^2} \Real\Big\{ \nabla_{\boldsymbol{\theta}_p} \vect{g}_p^H \cdot \mat{R}_X \cdot \nabla_{\boldsymbol{\theta}_p}^H \vect{g}_p \Big\} \tag{11}
\end{equation}
where $\mat{R}_X = \sum_k \vect{w}_k \vect{w}_k^H + \sum_p \vect{z}_p \vect{z}_p^H = \sum_k \mat{W}_k + \mat{Z}$ is the transmit covariance matrix and $\boldsymbol{\theta}_p$ denotes the target state parameters.

\begin{assumption}[Sensing Parameter Estimation Model]\label{ass:1}
This paper considers target state parameters $\boldsymbol{\theta}_p = [\theta_p]$ (single-angle estimation). Under this setting, the gradient matrix $\nabla_{\boldsymbol{\theta}_p} \vect{g}_p$ is determined by the target prediction state within the current time slot and can be treated as a known constant matrix. Consequently, $\mat{J}_p^{\text{data}}$ is an \textbf{affine function} of $\mat{R}_X$, and its trace constraint can be precisely formulated as:
\begin{equation}
\tr(\mat{F}_p \mat{R}_X) \geq \Gamma_{\text{Track},p}
\end{equation}
where $\mat{F}_p = \frac{2}{\sigma_s^2} \Real\Big\{ \nabla_{\boldsymbol{\theta}_p}^H \vect{g}_p \nabla_{\boldsymbol{\theta}_p} \vect{g}_p^H \Big\}$ is a known constant positive semidefinite matrix.
\end{assumption}

\begin{remark}
Since this work focuses on direction-of-arrival (DOA) estimation for targets, the relative delay and Doppler shift between the probing signal and the target are regarded as slowly varying parameters within a short observation frame. Under this condition, the FIM degenerates into a constant matrix $\mat{F}_p$ that depends solely on the angle gradient, thereby guaranteeing the convex affine property of the constraint.
\end{remark}

\textbf{Proof of FIM Linearity}: Let $\mat{G}_p = \nabla_{\boldsymbol{\theta}_p} \vect{g}_p^H \in \mathbb{C}^{D \times MN_t}$, then:
\begin{equation}
\mat{J}_p^{\text{data}} = \frac{2}{\sigma_s^2} \Real\{ \mat{G}_p \mat{R}_X \mat{G}_p^H \}
\end{equation}

Expanding the $(a,b)$-th element:
\begin{equation}
(\mat{J}_p^{\text{data}})_{ab} = \frac{2}{\sigma_s^2} \Real\Big\{ \sum_{i,j} (\mat{G}_p)_{ai} (\mat{R}_X)_{ij} (\mat{G}_p^H)_{jb} \Big\}
\end{equation}

This is a linear combination of the elements of $\mat{R}_X$. Therefore, the trace constraint:
\begin{equation}
\tr(\mat{J}_p^{\text{data}}) = \tr(\mat{F}_p \mat{R}_X) \tag{12}
\end{equation}

where:
\begin{equation}
\mat{F}_p = \frac{2}{\sigma_s^2} \Real\Big\{ \nabla_{\boldsymbol{\theta}_p}^H \vect{g}_p \nabla_{\boldsymbol{\theta}_p} \vect{g}_p^H \Big\} \tag{13}
\end{equation}

is a constant positive semidefinite matrix independent of the optimization variables. \hfill$\square$

%==============================================================================
\section{SDP Relaxation and Global Variable Reformulation (Plate II: Mathematical Reformulation)}
%==============================================================================

The physical modeling in Section II yields a highly non-convex problem (P1). This section presents a seven-step convexification chain that transforms (P1) into a standard convex SDP (P3).

%----- BEGIN Convexification chain insertion (EN) -----
%==============================================================================
% 6 Non-convex Sources Identification
%==============================================================================
\subsubsection{Identification of Non-convex Sources}

To make the subsequent convexification derivations amenable to rigorous review, this subsection pre-identifies all non-convex sources in the original problem and provides the mathematical justification for each. The original problem (P1) carries constraint labels (5a)--(5h); see \S\ref{sec:problem}.

\begin{table}[htbp]
\centering
\caption{Identification of non-convex sources in problem (P1)}\label{tab:nonconvex_sources}
\begin{tabular}{clll}
\toprule
ID & Non-convex source & Affected constraints & Convexity violation type \\
\midrule
NC1 & SINR fractional structure & (5b), (5c) & Convex set not closed under multiplication \\
NC2 & Semi-infinite robust constraints & (5b), (5c) & Infinite constraints violate convexity \\
NC3 & Binary combinatorial constraints & (5h) & Discrete set $\{0,1\}$ is non-convex \\
NC4 & Implicit rank-1 constraints & (5a)--(5f) & Rank-1 set is non-convex \\
NC5 & Beam-channel bilinear coupling & (5b), (5d) & Bilinear functions are non-convex \\
NC6 & Matrix inverse constraints & (5d) & Inverse map does not preserve convexity \\
\bottomrule
\end{tabular}
\end{table}

\textbf{Mathematical criterion for convexity}: A set $\mathcal{C}$ is convex iff $\forall \vect{x},\vect{y}\in\mathcal{C},\forall \theta\in[0,1]: \theta\vect{x}+(1-\theta)\vect{y}\in\mathcal{C}$. Using this criterion:

\begin{itemize}
    \item $\|\vect{w}\|_2^2 \leq \alpha$: convex (second-order cone)
    \item $\tr(\mat{H}_k \mat{W}_k) \geq \beta$ with $\mat{W}_k \succeq \mat{0}$: convex (PSD cone + linear function)
    \item $\rank(\mat{W}_k) = 1$: \textbf{non-convex} (rank-1 set is not affine)
    \item $\frac{x^2}{y} \leq \alpha, y > 0$: \textbf{non-convex} (fractional structure does not preserve convexity)
    \item $b \in \{0,1\}$: \textbf{non-convex} (discrete point set)
\end{itemize}

%==============================================================================
% 7-Step Convexification Chain
%==============================================================================
\subsubsection{Seven-Step Convexification Chain}

Targeting the six non-convex sources identified above, this subsection presents a numbered seven-step convexification chain. Each step specifies: which non-convex source it eliminates, the pre- and post-transformation mathematical forms, the transformation property (strictly equivalent / tight relaxation / conservative approximation), and a proposition with proof.

%------------------------------------------------------------------------------
\paragraph{Step 1: Global Variable Lifting (eliminate NC1, NC5)}
\textbf{Transformation type}: \textbf{Strictly equivalent}.

Lift the vector beam $\vect{w}_k$ into a covariance matrix:
\begin{equation}
\mat{W}_k = \vect{w}_k \vect{w}_k^H \in \mathbb{C}^{MN_t \times MN_t} \tag{15}
\end{equation}
\begin{equation}
\mat{Z} = \sum_p \vect{z}_p \vect{z}_p^H \in \mathbb{C}^{MN_t \times MN_t} \tag{16}
\end{equation}

\begin{proposition}[Lifting equivalence]\label{prop:lifting}
When $\mat{W}_k = \vect{w}_k \vect{w}_k^H$, we have $\tr(\mat{H}_k \mat{W}_k) = |\vect{h}_k^H \vect{w}_k|^2$, and the lifting $\vect{w}_k \mapsto \mat{W}_k$ is a one-to-one map.
\end{proposition}

\textbf{Convexity impact}: The objective $\sum_k \|\vect{w}_k\|^2 = \sum_k \tr(\mat{W}_k)$ becomes a linear function of $\mat{W}_k$, so the \textbf{objective becomes convex}.

%------------------------------------------------------------------------------
\paragraph{Step 2: SDR Relaxation (eliminate NC4)}
\textbf{Transformation type}: \textbf{Tight relaxation}. Strictly equivalent for $K\leq 2$; lower bound for $K>2$.

The rank-1 constraint $\mat{W}_k = \vect{w}_k \vect{w}_k^H$ is equivalent to
\begin{equation}
\mat{W}_k \succeq \mat{0}, \quad \rank(\mat{W}_k) = 1
\end{equation}
SDR drops the rank-1 constraint, requiring only
\begin{equation}
\mat{W}_k \succeq \mat{0} \tag{S2.1}
\end{equation}
The feasible set enlarges from $\mathcal{F}_{\text{rank-1}} = \{\mat{W}_k \succeq \mat{0} : \rank(\mat{W}_k)=1\}$ to $\mathcal{F}_{\text{SDR}} = \{\mat{W}_k \succeq \mat{0}\}$ (the convex PSD cone).

\begin{proposition}[SDR tightness conditions]\label{prop:sdr_tight}
SDR is equivalent to the original problem iff the optimal solution satisfies $\rank(\mat{W}_k^*) = 1$.
\begin{enumerate}
    \item SDR is tight for $K \leq 2$ (guaranteed by Sidiropoulos, Davidson, Luo 2006 *IEEE Trans. SP* Theorem 1 for MISO multicast).
    \item SDR is approximately tight in the high-SNR regime.
    \item For general $K>2$: Gaussian randomization ($\xi_l \sim \mathcal{CN}(\mat{0}, \mat{W}_k^*)$) extracts a feasible rank-1 solution with $O(1/L)$ performance loss.
\end{enumerate}
\end{proposition}

\textbf{Equivalence/lower/upper bound}: strictly equivalent when tight; otherwise a \textbf{lower bound} on the objective (enlarged feasible set $\Rightarrow$ optimal value $\leq$ original).

%------------------------------------------------------------------------------
\paragraph{Step 3: S-Procedure for Robust SINR (eliminate NC2)}
\textbf{Transformation type}: \textbf{Conservative approximation}. Trades a 1.75 dB power margin for a closed-form LMI.

The worst-case SINR constraint is a semi-infinite optimization:
\begin{equation}
\min_{\norm{\Delta\vect{h}_k} \leq \epsilon_h \norm{\hat{\vect{h}}_k}} \frac{|(\hat{\vect{h}}_k + \Delta\vect{h}_k)^H \mat{W}_k (\hat{\vect{h}}_k + \Delta\vect{h}_k)|}{\sum_{j\neq k}|(\hat{\vect{h}}_k + \Delta\vect{h}_k)^H \mat{W}_j(\hat{\vect{h}}_k + \Delta\vect{h}_k)| + \sigma_c^2} \geq \gamma_k
\end{equation}

\begin{lemma}[Cauchy-Schwarz worst case]\label{lem:cauchy_wc}
For any $\vect{x}, \vect{y}$ and $\Delta\vect{y}$ with $\|\Delta\vect{y}\| \leq \epsilon\|\vect{y}\|$:
\begin{equation}
\min_{\|\Delta\vect{y}\|\leq\epsilon\|\vect{y}\|} |\vect{x}^H(\vect{y}+\Delta\vect{y})|^2 = |\vect{x}^H\vect{y}|^2 (1-\epsilon)^2
\end{equation}
\textit{Proof}: By $| \vect{x}^H \Delta\vect{y}| \leq \|\vect{x}\| \cdot \|\Delta\vect{y}\| \leq \|\vect{x}\| \cdot \epsilon\|\vect{y}\|$, with equality at $\Delta\vect{y} = -\epsilon \frac{\|\vect{y}\|}{\|\vect{x}\|} \vect{x}$. $\square$
\end{lemma}

Applying Lemma~\ref{lem:cauchy_wc} (lower-bounding the numerator, upper-bounding the denominator) and cross-multiplying, the worst-case robust constraint becomes a linear constraint on a robust threshold $\gamma_k^{\text{robust}}$ (see Eq.~\eqref{eq:lmi_comm}).

\textbf{Equivalence/lower/upper bound}: conservative upper bound. The true worst-case SINR is higher than this approximation $\Rightarrow$ the solved feasible set is slightly larger than the true robust feasible set, but robustness is guaranteed under true channel realizations. The cost is approximately 1.75 dB power margin (for $\epsilon_h=0.10$, $\eta_h \approx 0.669$).

%------------------------------------------------------------------------------
\paragraph{Step 4: SINR Fraction Linearization (eliminate NC1, communication part)}
\textbf{Transformation type}: \textbf{Strictly equivalent} (provided the denominator $> 0$).

After Step 3's robust approximation, the SINR is still fractional. Cross-multiplication yields a linear matrix inequality:
\begin{equation}
\tr(\hat{\mat{H}}_k \mat{W}_k) - \gamma_k^{\text{robust}} \sum_{j\neq k} \tr(\hat{\mat{H}}_k \mat{W}_j) \geq \gamma_k^{\text{robust}} \sigma_c^2 \tag{S4.1}
\end{equation}
where $\hat{\mat{H}}_k = \hat{\vect{h}}_k \hat{\vect{h}}_k^H$.

\begin{proposition}[Fraction linearization equivalence]\label{prop:frac}
When the denominator $\sum_{j\neq k}|\hat{\vect{h}}_k^H \mat{W}_j \hat{\vect{h}}_k| + \sigma_c^2 > 0$,
\begin{equation}
\frac{A}{B} \geq \gamma \iff A \geq \gamma B
\end{equation}
\end{proposition}

\textbf{Convexity impact}: Eq.~\eqref{eq:lmi_comm} (equivalent to (S4.1) after S-Procedure lifting) becomes an LMI in $\mat{W}_k$, so the \textbf{communication constraint becomes convex}.

%------------------------------------------------------------------------------
\paragraph{Step 5a: PCRB Affine Expansion (eliminate NC6)}
\textbf{Transformation type}: \textbf{Strictly equivalent} (under Assumption 1).

The Fisher information matrix is an affine function of $\mat{R}_X$. Define
\begin{equation}
\mat{F}_p = \frac{2}{\sigma_s^2} \Real\Big\{ \nabla_{\boldsymbol{\theta}_p}^H \vect{g}_p \cdot \nabla_{\boldsymbol{\theta}_p} \vect{g}_p^H \Big\} \in \mathbb{S}_+^{D\times D} \tag{13}
\end{equation}

The PCRB constraint becomes
\begin{equation}
\tr(\mat{F}_p \mat{R}_X) \geq \Gamma_{\text{Track},p} \tag{S5.1}
\end{equation}

\begin{proposition}[PCRB linearization equivalence]\label{prop:pcrb}
Under Assumption 1 ($\nabla_{\boldsymbol{\theta}_p}\vect{g}_p$ is known within the current time slot), $\tr(\mat{J}_p^{\text{data}}) = \tr(\mat{F}_p \mat{R}_X)$.
\end{proposition}

\textbf{Convexity impact}: Eq.~(S5.1) is a linear constraint in $\mat{R}_X$, so the \textbf{sensing PCRB constraint becomes convex}.

%------------------------------------------------------------------------------
\paragraph{Step 5b: Sensing SINR Linearization}
\textbf{Transformation type}: \textbf{Strictly equivalent} (rank-1 optimal solution is matched filtering).

In the sensing SINR constraint $\tr(\vect{g}_p\vect{g}_p^H \mat{Z}) \geq \gamma_S^{\text{PoD}} \sigma_s^2$, the matrix $\vect{g}_p\vect{g}_p^H$ is a known constant PSD matrix and $\tr(\cdot)$ is linear in $\mat{Z}$.

\begin{proposition}[Sensing SINR linearization equivalence]\label{prop:sens_sinr}
When $\mat{Z} = \vect{z}\vect{z}^H$ (rank-1), the optimal sensing beam is $\vect{z}^* = \sqrt{P_S} \vect{g}_p/\|\vect{g}_p\|$ (matched filtering), and the constraint holds with equality.
\end{proposition}

%------------------------------------------------------------------------------
\paragraph{Step 6: Power Constraint Is Already Convex}
The power constraint $\sum_k \tr(\mat{W}_k) + \tr(\mat{Z}_m) \leq P_{\max}$ is linear in $\mat{W}_k, \mat{Z}_m$, \textbf{already convex}; no transformation needed.

%------------------------------------------------------------------------------
\paragraph{Step 7: AP Selection Two-Step Decomposition (handle NC3)}
\textbf{Transformation type}: \textbf{Heuristic lower bound}. The outer discrete AP assignment is NP-hard; the inner continuous problem is a convex SDP for any fixed AP set.

This work adopts a two-step decomposition:
\begin{enumerate}
    \item \textbf{Outer (discrete)}: sort APs by large-scale fading $\text{PL}(d_{m,p})$ and select the top-$N_{\text{req}}$ APs to fix $\mathcal{M}^{\text{all}}$.
    \item \textbf{Inner (convex)}: solve the convex SDP (P3) on the fixed AP set.
\end{enumerate}

\begin{remark}[Engineering trade-off for AP selection]\label{rem:ap_heuristic}
The AP selection problem is NP-hard (a variant of the $K$-medoid problem). This work uses an outer heuristic plus inner convex SDP to obtain a polynomial-complexity, engineering-acceptable sub-optimal solution. The complexity lower bound is established in the companion document~\cite{convex_chain} \S6.
\end{remark}

%==============================================================================
% Convexification Chain Summary Table
%==============================================================================
\subsubsection{Convexification Chain Summary}

The following table summarizes the transformation type and performance impact of each step:

\begin{table}[htbp]
\centering
\caption{Summary of the seven-step convexification chain}\label{tab:convex_chain}
\begin{tabular}{clll}
\toprule
Step & Transformation & Eliminates & Type \\
\midrule
1 & Global variable lifting $\vect{w}\vect{w}^H \to \mat{W}$ & NC1, NC5 & Strictly equivalent \\
2 & SDR relaxation (drop rank-1) & NC4 & Tight relaxation (equiv.\ for $K\leq 2$) \\
3 & S-Procedure closed-form bound & NC2 & Conservative approx.\ (1.75 dB margin) \\
4 & SINR fraction linearization & NC1 (comm.) & Strictly equivalent \\
5a & PCRB affine expansion & NC6 & Strictly equivalent (Assumption 1) \\
5b & Sensing SINR linearization & NC1 (sens.) & Strictly equivalent \\
6 & Power constraint (already convex) & --- & Already convex \\
7 & AP selection two-step decomp. & NC3 & Heuristic lower bound \\
\bottomrule
\end{tabular}
\end{table}

\textbf{Key conclusion}: The physical-layer convexification (Steps 1--6) is almost entirely strictly equivalent. Only Step 3 is a conservative engineering approximation, and Step 2 is tight for $K \leq 2$. AP selection (Step 7) uses a heuristic to maintain polynomial complexity.

%==============================================================================
% Final Convex SDP Problem (P3)
%==============================================================================
\subsubsection{Final Convex SDP Problem (P3)}

After the seven-step chain, the original problem (P1) reduces to the following \textbf{standard convex SDP}:

\begin{align}
\min_{\{\mat{W}_k\}, \mat{Z}, \{\mu_k\}} \quad & \sum_{k=1}^K \tr(\mat{W}_k) + \tr(\mat{Z}) \tag{P3} \\
\text{s.t.} \quad & \begin{bmatrix} \mat{A}_k + \mu_k \mat{I} & \mat{A}_k \hat{\vect{h}}_k \\[1.5ex] \hat{\vect{h}}_k^H \mat{A}_k & \hat{\vect{h}}_k^H \mat{A}_k \hat{\vect{h}}_k - \sigma_c^2 - \mu_k \epsilon_h^2 \end{bmatrix} \succeq \mat{0}, \quad \forall k \label{eq:lmi_comm} \\
& \tr(\vect{g}_p \vect{g}_p^H \mat{Z}) \geq \gamma_S^{\text{PoD}} \sigma_s^2, \quad \forall p \label{eq:sinr_sens_sdp} \\
& \tr(\mat{F}_p \mat{R}_X) \geq \Gamma_{\text{Track},p}, \quad \forall p \label{eq:pcrb_sdp} \\
& \tr(\mat{E}_m \mat{R}_X) \leq P_{\max}, \quad \forall m \label{eq:power_sdp} \\
& \mat{W}_k \succeq \mat{0}, \quad \forall k \label{eq:psd_w} \\
& \mat{Z} \succeq \mat{0} \label{eq:psd_z} \\
& \mu_k \geq 0, \quad \forall k \label{eq:mu_nonneg}
\end{align}

where $\mat{A}_k = \frac{1}{\gamma_k} \mat{W}_k - \sum_{j \neq k} \mat{W}_j$, $\mat{F}_p$ is given by Eq.~(13), $\mat{E}_m$ is the selection matrix for AP $m$, and $\mu_k \geq 0$ is the S-Procedure slack variable.

\begin{table}[htbp]
\centering
\caption{Problem (P3) constraint structure and variable mapping}\label{tab:p3_structure}
\begin{tabular}{clll}
\toprule
Constraint & Type & Variables & Physical meaning \\
\midrule
\eqref{eq:lmi_comm} & LMI ($MN_t+1$ dim.) & $\mat{W}_k, \mu_k$ & Communication worst-case robustness \\
\eqref{eq:sinr_sens_sdp} & Linear inequality & $\mat{Z}$ & Sensing detection probability \\
\eqref{eq:pcrb_sdp} & Linear inequality & $\mat{R}_X$ & Tracking accuracy (PCRB) \\
\eqref{eq:power_sdp} & Linear inequality & $\mat{R}_X$ & Per-AP peak power \\
\eqref{eq:psd_w}--\eqref{eq:psd_z} & PSD cone & $\mat{W}_k, \mat{Z}$ & SDR relaxation \\
\eqref{eq:mu_nonneg} & Non-negative orthant & $\mu_k$ & S-Procedure slack \\
\bottomrule
\end{tabular}
\end{table}

\begin{theorem}[Convexity of the problem]\label{thm:convex_p3}
Problem (P3) is a standard convex SDP: the objective is linear and all constraints are linear matrix inequalities (LMIs) or semidefinite cone constraints. It can be solved to arbitrary accuracy in polynomial time by standard SDP solvers (MOSEK / SeDuMi / SDPT3). The solution complexity is $O((MN_t)^{3.5})$, taking approximately 5--10 seconds for $M=16, N_t=4$.
\end{theorem}

\begin{theorem}[SDR tightness]\label{thm:sdr_tight}
For the total power minimization problem, SDR is tight when $K \leq 2$, i.e., the optimal solution satisfies $\rank(\mat{W}_k^*) = 1$. In the high-SNR regime, SDR is approximately tight. For general $K>2$, beam recovery via Gaussian randomization incurs a performance loss upper bounded by $O(1/L)$ (where $L$ is the number of candidates).
\end{theorem}

%==============================================================================
% Complexity Comparison
%==============================================================================
\subsubsection{Complexity Comparison: Before and After Convexification}

\begin{table}[htbp]
\centering
\caption{Comparison of problem form and solution complexity}\label{tab:complexity_compare}
\begin{tabular}{llll}
\toprule
Problem & Form & Solver & Complexity \\
\midrule
(P1) Original & MINLP (mixed-integer non-convex) & No general algorithm & NP-hard \\
(P3) Convex SDP & Convex semidefinite program & MOSEK / SeDuMi & $O((MN_t)^{3.5})$ \\
\bottomrule
\end{tabular}
\end{table}

Convexification transforms an NP-hard mixed-integer non-convex problem into a polynomial-time-solvable standard convex SDP, which is the central value of the convexification methodology.

%----- END convexification chain insertion -----

\subsection{From Beams to Covariances}

To transform the non-convex beamforming problem into a convex optimization problem, we introduce covariance matrix variables:

\textbf{Communication covariance matrix}:
\begin{equation}
\mat{W}_k = \vect{w}_k \vect{w}_k^H \in \mathbb{C}^{MN_t \times MN_t} \tag{15}
\end{equation}

\textbf{Sensing covariance matrix}:
\begin{equation}
\mat{Z} = \sum_p \vect{z}_p \vect{z}_p^H \in \mathbb{C}^{MN_t \times MN_t} \tag{16}
\end{equation}

\textbf{Global transmit covariance}:
\begin{equation}
\mat{R}_X = \sum_{k=1}^K \mat{W}_k + \mat{Z} \tag{17}
\end{equation}

\subsection{SDR Relaxation and Final Convex SDP Problem}

The original problem requires $\rank(\mat{W}_k) = 1$ (since $\mat{W}_k = \vect{w}_k \vect{w}_k^H$). Semidefinite relaxation (SDR) discards the rank-one constraint, requiring only $\mat{W}_k \succeq \mat{0}$. Furthermore, through the S-Procedure, the infinite-dimensional worst-case SINR constraint is exactly transformed into a finite-dimensional linear matrix inequality (LMI), yielding the following \textbf{standard convex semidefinite programming problem} that can be directly fed into CVX/MOSEK.

\textbf{Core Auxiliary Definition}:
\begin{equation}
\mat{A}_k = \frac{1}{\gamma_k} \mat{W}_k - \sum_{j \neq k} \mat{W}_j
\end{equation}

\textbf{Final Convex SDP Problem (P3)}:
\begin{align}
\min_{\{\mat{W}_k\}, \mat{Z}, \{\mu_k\}} \quad & \sum_{k=1}^K \tr(\mat{W}_k) + \tr(\mat{Z}) \tag{P3} \\
\text{s.t.} \quad & \begin{bmatrix} \mat{A}_k + \mu_k \mat{I} & \mat{A}_k \hat{\vect{h}}_k \\[1.5ex] \hat{\vect{h}}_k^H \mat{A}_k & \hat{\vect{h}}_k^H \mat{A}_k \hat{\vect{h}}_k - \sigma_c^2 - \mu_k \epsilon_h^2 \end{bmatrix} \succeq \mat{0}, \quad \forall k \label{eq:lmi_comm} \\
& \tr(\vect{g}_p \vect{g}_p^H \mat{Z}) \geq \gamma_S^{\text{PoD}} \sigma_s^2, \quad \forall p \label{eq:sinr_sens_sdp} \\
& \tr(\mat{F}_p \mat{R}_X) \geq \Gamma_{\text{Track},p}, \quad \forall p \label{eq:pcrb_sdp} \\
& \tr(\mat{E}_m \mat{R}_X) \leq P_{\max}, \quad \forall m \label{eq:power_sdp} \\
& \mat{W}_k \succeq \mat{0}, \quad \forall k \label{eq:psd_w} \\
& \mat{Z} \succeq \mat{0} \label{eq:psd_z} \\
& \mu_k \geq 0, \quad \forall k \label{eq:mu_nonneg}
\end{align}

where $\mat{E}_m$ is the selection matrix for AP $m$, and $\mu_k \geq 0$ is the S-Procedure slack variable.

\begin{table}[htbp]
\centering
\caption{Problem (P3) Constraint Structure and Variable Mapping}\label{tab:p3_structure}
\begin{tabular}{clll}
\toprule
Constraint & Type & Variables & Physical Meaning \\
\midrule
\eqref{eq:lmi_comm} & LMI ($MN_t+1$ dim.) & $\mat{W}_k, \mu_k$ & Communication worst-case robustness \\
\eqref{eq:sinr_sens_sdp} & Linear inequality & $\mat{Z}$ & Sensing detection probability \\
\eqref{eq:pcrb_sdp} & Linear inequality & $\mat{R}_X$ & Tracking accuracy (PCRB) \\
\eqref{eq:power_sdp} & Linear inequality & $\mat{R}_X$ & Per-AP peak power \\
\eqref{eq:psd_w}--\eqref{eq:psd_z} & PSD cone & $\mat{W}_k, \mat{Z}$ & SDR relaxation \\
\eqref{eq:mu_nonneg} & Non-negative orthant & $\mu_k$ & S-Procedure slack \\
\bottomrule
\end{tabular}
\end{table}



\subsection{Tightness Conditions}

\begin{theorem}[SDR Tightness]\label{thm:1}
For the total power minimization problem, SDR provides a lower bound $P_{\text{SDR}}^* \leq P_{\text{original}}^*$. No universal SDR tightness theorem applies to our per-user SINR + cell-free cooperative + robust problem—the multicast $G_k \leq 2$ tightness of Karipidis, Sidiropoulos, Luo 2008 and the robust MISO $N_t \leq 2$ tightness of Tuan, Apaydin, Lu 2012 (*Eurasip JWCN*) do not transfer.
\end{theorem}

\begin{remark}
In the high-SNR regime, SDR is approximately tight with controllable performance loss. For the general case, beam recovery via Gaussian randomization (the standard practice in SDR-based beamforming, see Luo, Ma, So, Ye, Zhang 2010 *IEEE Signal Process. Mag.*) recovers a feasible rank-one candidate with **no tight performance-loss bound**; engineering practice uses $L = 100 \sim 1000$ candidates and selects the best.
\end{remark}

%==============================================================================
\section{S-Procedure Robust Constraint Transformation}
%==============================================================================

\subsection{Communication Robust Constraint (Exact Transformation)}

The worst-case SINR constraint is:
\begin{equation}
\min_{\norm{\Delta\vect{h}_k} \leq \epsilon_h} \frac{|(\hat{\vect{h}}_k + \Delta\vect{h}_k)^H \vect{w}_k|^2}{\sum_{j \neq k} |(\hat{\vect{h}}_k + \Delta\vect{h}_k)^H \vect{w}_j|^2 + \sigma_c^2} \geq \gamma_k
\end{equation}

Using the S-Procedure, this infinite-dimensional constraint can be transformed into a finite-dimensional LMI. Define:
\begin{equation}
\mat{A}_k = \frac{1}{\gamma_k} \mat{W}_k - \sum_{j \neq k} \mat{W}_j
\end{equation}

Then there exists $\mu_k \geq 0$ such that the following LMI holds:
\begin{equation}
\begin{bmatrix} \mat{A}_k + \mu_k \mat{I} & \mat{A}_k \hat{\vect{h}}_k \\[1.5ex] \hat{\vect{h}}_k^H \mat{A}_k & \hat{\vect{h}}_k^H \mat{A}_k \hat{\vect{h}}_k - \sigma_c^2 - \mu_k \epsilon_h^2 \end{bmatrix} \succeq \mat{0} \tag{18}
\end{equation}

\begin{proposition}
The S-Procedure is an \textbf{exactly equivalent} transformation, not an approximation. That is, the original worst-case constraint and the LMI constraint are completely equivalent.
\end{proposition}

\subsection{Sensing Robust Constraint (Optional)}

Similarly, for the sensing channel, the LMI form is:
\begin{equation}
\begin{bmatrix} \frac{1}{\gamma_S^{\text{PoD}}} \mat{Z} + \nu_p \mat{I} & \frac{1}{\gamma_S^{\text{PoD}}} \mat{Z} \hat{\vect{g}}_p \\[1.5ex] \frac{1}{\gamma_S^{\text{PoD}}} \hat{\vect{g}}_p^H \mat{Z} & \frac{1}{\gamma_S^{\text{PoD}}} \hat{\vect{g}}_p^H \mat{Z} \hat{\vect{g}}_p - \sigma_s^2 - \nu_p \epsilon_g^2 \end{bmatrix} \succeq \mat{0} \tag{19}
\end{equation}

where $\nu_p \geq 0$ is the sensing S-Procedure slack variable.

\begin{remark}[LMI Typesetting Optimization]
The matrix elements in (19) contain denominators $\gamma_S^{\text{PoD}}$. In standard LMI notation, both sides of the inequality can be multiplied by $\gamma_S^{\text{PoD}}$, yielding $\mat{Z} + \nu_p \gamma_S^{\text{PoD}} \mat{I}$ in the upper-left corner, thereby avoiding fractional structures inside the matrix for more elegant typesetting. Mathematical equivalence is preserved.
\end{remark}

\begin{assumption}[Sensing Channel Perfection]\label{ass:2}
This work focuses on robust design for the communication link, assuming that sensing channels are self-calibrated through tracking echoes within the current time slot with negligible error. Specifically:
\begin{enumerate}
    \item Target state parameters $\boldsymbol{\theta}_p$ are accurately predicted within a short time slot, with prediction errors much smaller than one wavelength
    \item Sensing channels $\vect{g}_p$ are calibrated through tracking echoes from the previous time slot, with errors incorporated into the next slot update
    \item Sensing tasks adopt a "detection-tracking" cascade architecture: conservative thresholds during detection, and time-slot smoothness compensates for errors during tracking
\end{enumerate}
\end{assumption}

\begin{remark}
Sensing channels are self-calibrating (the probing signal is transmitted and the echo is received, with the echo carrying current channel state), whereas communication channels rely on uplink pilot estimation, where pilot contamination leads to error accumulation. In top-tier journal system modeling, focusing on communication-side pilot contamination while assuming the radar side benefits from LOS echoes and Kalman tracking filtering for precise prediction is a common and prudent compromise.
\end{remark}

%==============================================================================
\section{Proof of Sensing Constraint Convexity}
%==============================================================================

This section proves the convexity of the sensing-side constraints.

\subsection{PCRB Constraint Convexity}

\textbf{Constraint}: $\tr(\mat{F}_p \mat{R}_X) \geq \Gamma_{\text{Track},p}$

\begin{itemize}
    \item $\mat{F}_p$ is a known constant positive semidefinite matrix
    \item $\mat{R}_X$ is the optimization variable (positive semidefinite matrix)
    \item $\tr(\mat{F}_p \mat{R}_X)$ is a linear function of $\mat{R}_X$
\end{itemize}

\textbf{Conclusion}: Linear inequality constraint, naturally convex.

\subsection{Sensing SINR Constraint Convexity}

\textbf{Constraint}: $\tr(\vect{g}_p \vect{g}_p^H \mat{Z}) \geq \gamma_S^{\text{PoD}} \sigma_s^2$

\begin{itemize}
    \item $\vect{g}_p \vect{g}_p^H$ is a known constant positive semidefinite matrix
    \item $\mat{Z}$ is the optimization variable
    \item $\tr(\vect{g}_p \vect{g}_p^H \mat{Z})$ is a linear function of $\mat{Z}$
\end{itemize}

\textbf{Conclusion}: Linear inequality constraint, naturally convex.

\subsection{Power Constraint Convexity}

\textbf{Constraint}: $\tr(\mat{E}_m \mat{R}_X) \leq P_{\max}$

\begin{itemize}
    \item $\mat{E}_m$ is a known constant matrix
    \item $\mat{R}_X$ is the optimization variable
\end{itemize}

\textbf{Conclusion}: Linear inequality constraint, naturally convex.

\subsection{Summary}

Table~\ref{tab:convexity} summarizes the convexity of each constraint.

\begin{table}[htbp]
\centering
\caption{Convexity Summary of Constraints}\label{tab:convexity}
\begin{tabular}{lll}
\toprule
Constraint & Form & Convexity \\
\midrule
Communication Robust (S-Procedure) & LMI & Convex \\
PCRB & Linear trace & Convex \\
Sensing SINR & Linear trace & Convex \\
Power & Linear trace & Convex \\
Positive Semidefinite & $\mat{W}_k, \mat{Z} \succeq \mat{0}$ & Convex \\
\bottomrule
\end{tabular}
\end{table}

\begin{theorem}[Final Problem Convexity]
Problem (P3) is a standard convex SDP, solvable to arbitrary precision via interior-point methods in polynomial time. Combined with Algorithm~\ref{alg:rank1} (Gaussian randomization rank-one recovery), this theoretical framework can be directly translated into MATLAB/CVX code.
\end{theorem}

%==============================================================================


\end{document}
